\chapter*{前言}
\addcontentsline{toc}{chapter}{前言}
本书的目的在于呈现一大批计算刚体动力学的最高效算法，并对它们作足够详细的讲解，使读者既能理解这些算法的工作原理，又能懂得如何改造它们（或创建新算法）以满足自身的需要。书中收录了以下几个著名算法：递推牛顿–欧拉算法、复合刚体算法和关节体算法。此外还包括处理运动环和浮动基座的算法。每个算法都从第一性原理出发推导得到，并且既以方程组的形式给出，也以伪代码程序的形式给出；后者经过精心设计，便于翻译成任何合适的编程语言。本书还讲解了建立刚体系统运动方程所用的一些数学技术。特别是，本书展示了如何用六维（6D）向量来表述动力学，并阐释了作为最高效算法之基础的递推形式。其他主题包括：如何构建刚体系统的计算机模型；如何利用惯量矩阵中的稀疏性；关节体惯量这一概念；动力学计算中舍入误差的来源；以及刚体之间物理接触与碰撞的动力学。

刚体动力学往往容易演变成代数推导的海洋。不过，这主要是使用三维向量所致，改用六维向量记法即可在很大程度上克服这一弊病。本书采用基于空间向量的记法，将刚体运动的线性与角度两个方面统一到同一组量和方程之中。其结果通常是代数推导量缩减为原来的四分之一到六分之一。这种好处同样体现在计算机代码上：程序更短、更清晰，更易于阅读、编写和调试，而效率仍与使用标准三维向量的代码相当。

本书面向广泛的读者群体，从高年级本科生到研究人员和专业人士均可阅读。假定读者已具备一些刚体动力学的先备知识，例如通过一门动力学导论课程获得的知识，或通过阅读某本入门教材前几章获得的知识。不过，读者无需预先掌握六维向量的知识，因为这一主题将在第 2 章中从零讲起。本书也确实包含一些较为高深的内容，或许会引起动力学

专家以及六维向量学者的兴趣。本书不随附任何软件，但读者可以从作者网站上获取本书所述大多数算法的源代码。本书最初本打算作为《Robot Dynamics Algorithms》一书的第二版，该书早在 1987 年就已出版；但很快便显而易见，新素材已足够丰富，足以支撑撰写一本全新的著作。与其前身相比，本书中最值得注意的新内容包括：算法的显式伪代码描述；讲述如何对刚体系统建模的一章；利用分支诱导稀疏性的算法；对运动环和浮动基座系统的扩充论述；平面向量（空间向量的平面版本）；数值误差与模型灵敏度；以及关于如何实现空间向量运算的指导。
